% Authoritative LaTeX manuscript; Markdown is a generated compatibility projection.

% <!-- part-title: Modeling -->

% <!-- chapter-title: Behavioral Models: What May Happen? -->

\chapter{Behavioral Models: What May Happen?}

A structural model describes relationships among parts of a system. A behavioral model describes how the system acts or changes over time. Depending on the question, we might model messages exchanged in a protocol, steps in a workflow, or a quantity that changes continuously. Here we focus on a common choice for software systems: discrete states and the events that move a system between them. With such a model, we can ask whether a job can finish before it starts, or whether a job that starts will ever finish.

\section{States, Transitions, and Behavioral Analysis}\label{heading:states-transitions-and-behavioral-analysis}

Suppose a service processes jobs that may be waiting, running, completed, or failed. Can a job report completion before anyone has started it? To find out, we need to know where a job begins, which event starts it, and from which states the service permits completion. We can set aside the database schema, scheduling algorithm, and retry machinery for now. If a job can finish too early, we ought to be able to see how in a model small enough to trace by hand.

\subsection{Constructing a Behavioral Model}\label{heading:constructing-a-behavioral-model}

% <!-- point: the-reduction-is-what-makes-behavior-checkable | Reducing a service to states and transitions turns "can it reach a state it should never occupy?" from a reconstruction over scattered handlers into a question a person or tool can answer. -->

Four states and three events suffice. Let \(S=\{\mathrm{IDLE},\mathrm{RUNNING},\mathrm{DONE},\mathrm{FAILED}\}\) be the states, with each job starting in \(\mathrm{IDLE}\). The events are \(E=\{\mathrm{start},\mathrm{finish},\mathrm{fail}\}\). We still need to say which events are permitted in each state and where they take the job.

% <!-- point: transition-semantics-can-construct-intended-behavior | A partial transition function defines the effects of start, finish, and fail in a four-state job lifecycle. Undefined pairs admit no transition, and the same function can guide implementation, interpretation, or test generation. | terms: semantic-commitment, modeling-principle -->

What an event does depends on the state in which it occurs. The following partial function \(\delta:S\times E\rightharpoonup S\) gives the next state for each permitted pair:
\[
\begin{aligned}
\delta(\mathrm{IDLE},\mathrm{start}) &= \mathrm{RUNNING},\\
\delta(\mathrm{RUNNING},\mathrm{finish}) &= \mathrm{DONE},\\
\delta(\mathrm{RUNNING},\mathrm{fail}) &= \mathrm{FAILED}.
\end{aligned}
\]
Try \(\mathrm{finish}\) from \(\mathrm{IDLE}\). There is no entry for that pair, so the model forbids the transition. Now try \(\mathrm{start}\) followed by \(\mathrm{finish}\). The first entry takes the job to \(\mathrm{RUNNING}\), from which the second takes it to \(\mathrm{DONE}\). Both \(\mathrm{DONE}\) and \(\mathrm{FAILED}\) are terminal.

The model settles when a job may finish while leaving the engineer free to choose how to implement it. An agent might store the state in a database, manage jobs through a queue, or use an in-memory scheduler. A state-machine interpreter can consult the entries directly. A test generator can use those same entries to produce a sequence of events and predict the state the implementation should report afterward.

A failed attempt need not mean a failed job. In \autoref{canonical-behavioral}, a repairable failure can move a running job into Retrying, from which a new attempt returns it to Running. A job can therefore succeed after an earlier attempt failed. The extra state earns its place when we want to reason about recovery.

\begin{figure}[htbp]
\centering
\includegraphics[width=\linewidth]{book/latex/assets/canonical-behavioral-c9de560202.pdf}
\caption{\emph{The behavioral move: states and transitions.} This example specifies an allowed job lifecycle: a job passes through Running before Succeeded, with retry and failure transitions describing other possible paths. The legend lists example properties to investigate.}
\label{canonical-behavioral}
\end{figure}

\subsection{Properties and Analyses}\label{heading:behavioral-properties-and-analyses}

% <!-- point: semantics-are-states-and-permitted-transitions-and-the-omissions-are-the-obligations | A behavioral model represents states and transitions; it can specify allowed behavior or describe behavior to be checked against a requirement. Safety and liveness properties ask different questions over the admitted executions. -->

Starting at Pending and following permitted arrows gives us one possible execution of the model. Trace Pending, Running, Retrying, Running, Succeeded. That path shows that success after a retry is \emph{reachable}. Here the arrows specify behavior the service permits. In a model derived from an implementation, the same path would describe behavior to evaluate against a requirement.

One path can establish reachability, and one bad path can disprove a \textbf{safety} requirement. Suppose the requirement is that a job must enter Running before it reaches Succeeded. Finding one execution that does so leaves all the other executions to consider. Instead, search for a path from Pending to Succeeded that avoids Running. Any such path would show exactly how premature completion could occur; in the figure, there is none.

A job can obey every ordering rule and still never finish. The requirement that every started job eventually succeeds or fails is a \textbf{liveness} property. A limit of three retries bounds the number of attempts, but says nothing about how long a worker may remain in Running. Observing one slow attempt cannot settle whether it will ever finish.

% <!-- point: behavior-has-a-long-repertoire-of-state-transition-representations | State-transition models make lifecycle and recovery questions explicit; state machines, labeled transition systems, and statecharts represent behavior, while temporal specifications state properties over executions. -->

To investigate progress, we must say what can keep a job waiting and what allows it to resume. If the service relies on a timeout, represent what happens when it fires and state the assumption that it will be delivered. We can then ask whether any execution under those assumptions can avoid both terminal states indefinitely. Temporal model checking evaluates such requirements over executions~\citep{baier-katoen2008}. For a more complex lifecycle, statecharts can organize related states hierarchically and represent concurrent activity~\citep{harel1987}.

\section{DocAble: Job Lifecycles and Recovery}\label{heading:docable-job-lifecycles-and-recovery}

DocAble distributes document accessibility work across workers that may fail or lose their assignments. Recovery raises two behavioral questions beyond the simple job lifecycle in \hyperref[2.4-behavioral:heading:states-transitions-and-behavioral-analysis]{the preceding subsection}. Can a worker publish a recovered result in an unsafe order? Can a former owner return and interfere with a newer assignment? The first requires a model of publication states; the second requires a model of ownership across competing workers.

\textbf{Publishing a recovered result}

A worker recovering a document job may have a corrected artifact available locally before uploading it to persistent storage. If the worker updates the database record first and then crashes, the record can point to an artifact that other workers cannot retrieve. What ordering of publication and database updates prevents this failure?

\textbf{Recovery-ordering model · Behavioral}

\begin{itemize}
\tightlist
\item
  \textbf{Engineering question} --- In what order may recovery publish an artifact and update the record that refers to it?
\end{itemize}



\begin{itemize}
\tightlist
\item
  \textbf{Model} --- a transition system over the recovery lifecycle; the decisive fact is a transition that does \emph{not} exist.
\end{itemize}



\begin{itemize}
\tightlist
\item
  \textbf{Property} --- the database record may refer to the fallback artifact only after that artifact has been uploaded to persistent storage.
\end{itemize}



\begin{itemize}
\tightlist
\item
  \textbf{Quality attribute} --- data integrity; crash-safety.
\end{itemize}

% <!-- point: the-decisive-fact-is-an-absent-edge-between-staged-and-durable | The recovery obligation lives in the gap between local staging and durable publication: there is no direct FALLBACK_STAGED to DB_UPDATED edge, so the model excludes executions in which the record points at a not-yet-durable artifact. -->

\autoref{recovery-ordering-model} represents the recovery operation using three states. In \texttt{FALLBACK\_STAGED}, the corrected artifact exists locally; in \texttt{FALLBACK\_UPLOADED}, it has been published to durable storage; and in \texttt{DB\_UPDATED}, the database refers to it. The permitted path passes through all three states. The prohibited transition from \texttt{FALLBACK\_STAGED} directly to \texttt{DB\_UPDATED} represents the unsafe ordering: the database would refer to an artifact that had not yet been published.

\begin{figure}[htbp]
\centering
\includegraphics[width=\linewidth]{book/latex/assets/recovery-ordering-model-70e2639fc3.pdf}
\caption{\emph{Recovery-ordering model.} A recovered artifact must be published to durable storage before the database may refer to it. The prohibited transition from \texttt{FALLBACK\_STAGED} directly to \texttt{DB\_UPDATED} illustrates the ordering violation; the property also excludes indirect paths that bypass \texttt{FALLBACK\_UPLOADED}.}
\label{recovery-ordering-model}
\end{figure}

To check the ordering property, search for an execution that reaches \texttt{DB\_UPDATED} without first visiting \texttt{FALLBACK\_UPLOADED}. The three-state model has no such execution. In a larger lifecycle, checking only for the direct forbidden transition would be insufficient: another sequence of transitions might reach \texttt{DB\_UPDATED} while bypassing publication. This is a safety property, not a liveness guarantee. A worker could wait or retry indefinitely without violating the ordering requirement.

% <!-- point: the-recovery-ordering-is-a-slice-of-a-larger-lifecycle | The recovery model keeps the states needed to check publication order; a fuller job lifecycle can support broader workflow questions when it includes the relevant transitions. -->

The three-state recovery model is a reduction for checking publication order. DocAble's fuller job lifecycle includes queued, leased, remediating, and completed or failed states. That model can address broader workflow questions when it includes the relevant transitions; the smaller recovery model keeps only what the ordering question needs.

\textbf{Distinguishing a current owner from a returning one}

% <!-- point: ownership-has-a-long-vocabulary-of-claims-and-leases | Concurrent workers can acquire, lose, and regain rights to a job; locks, leases, capabilities, and claims express different aspects of access and ownership, which the model selects according to the question. -->

DocAble leases each in-flight job to one worker. Suppose worker A stalls and its lease expires. Worker B acquires the job, but A later returns and attempts to release its old assignment. How can the system distinguish A's obsolete claim from B's current one? \autoref{ownership-constructs} distinguishes several constructs for representing access and ownership. DocAble combines two: a claim records the assignment, and a lease limits its lifetime.

{\footnotesize
\begin{longtable}[]{@{}
  >{\raggedright\arraybackslash}p{(\linewidth - 2\tabcolsep) * \real{0.2400}}
  >{\raggedright\arraybackslash}p{(\linewidth - 2\tabcolsep) * \real{0.7600}}@{}}
\caption[Constructs for access and ownership]{\emph{Constructs for access and ownership.} Locks, leases, capabilities, and claims express different aspects of a worker\textquotesingle s right to act.}\label{ownership-constructs}\tabularnewline
\toprule\noalign{}
\begin{minipage}[b]{\linewidth}\raggedright
Construct
\end{minipage} & \begin{minipage}[b]{\linewidth}\raggedright
What it expresses
\end{minipage} \\
\midrule\noalign{}
\endhead
\bottomrule\noalign{}
\endlastfoot
Lock & Coordinates access to a resource. An exclusive lock admits one holder at a time. \\
Lease & Grants a right for a limited period; expiry can make that right available for reassignment~\citep{graycheriton1989}. \\
Capability & Carries authority to perform specified operations on a resource; that authority need not be exclusive. \\
Claim & Records a worker's assignment to a job and the conditions under which that assignment is valid. \\
\end{longtable}
}

% <!-- point: leased-iff-owner-is-the-models-definition | A job is leased exactly when it has an owner; the two are the same fact from two sides, so the model represents ownership as the presence of a current claim — this is what "leased" means, and it sets up the invariant that follows. -->

The model represents ownership through a current claim relating a worker to a job. A job is leased exactly when it has such a claim; ownership and lease status are not separate facts that can disagree within this model. Each claim also records a generation and a lifetime. Expiry makes the job available for reassignment, and a new assignment advances the generation. The generation distinguishes successive claims even when they concern the same job.

\textbf{In-flight lease · Ownership}

\begin{itemize}
\tightlist
\item
  \textbf{Engineering question} --- Who currently owns this in-flight work, and can a returning stale owner disturb the current claim?
\end{itemize}



\begin{itemize}
\tightlist
\item
  \textbf{Model} --- a claim relating a worker to a job, stamped with a generation and a lifetime; reclaim advances the generation.
\end{itemize}



\begin{itemize}
\tightlist
\item
  \textbf{Property} --- a stale claim cannot clear or supersede the current one.
\end{itemize}



\begin{itemize}
\tightlist
\item
  \textbf{Quality attribute} --- correctness under concurrency; crash-safety.
\end{itemize}

% <!-- point: the-epoch-generation-makes-a-stale-release-inapplicable | Suppose A holds generation 7, stalls, and is reclaimed; the job is re-leased to B at generation 8; a late release from A carrying generation 7 must not clear B's generation-8 claim, and the generation makes that decidable. | terms: soft-vs-hard-governance -->

\autoref{in-flight-ownership-states} follows one possible execution. Worker A acquires generation 7, then stalls until its lease expires. Worker B reclaims the job at generation 8. When A returns and attempts to release generation 7, the operation must not clear B's current claim. The generation number makes the stale operation distinguishable from a release by the current owner. Without that distinction, A's delayed release could clear B's valid assignment.

\begin{figure}[htbp]
\centering
\includegraphics[width=\linewidth]{book/latex/assets/in-flight-ownership-states-01b169485d.pdf}
\caption{\emph{Ownership under contention.} Worker A's generation-7 claim expires and worker B acquires generation 8. A later release from A cannot clear B's current claim. Recording the generation makes stale operations distinguishable from operations authorized by the current owner.}
\label{in-flight-ownership-states}
\end{figure}

% <!-- point: a-claim-carries-the-facts-that-make-contention-decidable | A claim binds a holder to a resource and carries which holder, which generation of claim, and when it expires — the generation being what tells a current claim from a superseded one — and because the hard failures live in interleavings, the characteristic analysis is state-space exploration with invariant checking. -->

The illustrated execution is only one possible ordering of events. To investigate the invariant, explore the interleavings of lease expiry, reassignment, and delayed release. At each reachable configuration, check whether a release from an earlier generation can clear or supersede the current claim. A counterexample would identify an event sequence that violates ownership; checking the reachable configurations establishes the property within the model's declared scope.

The publication model separates a local artifact from a durably stored one. The ownership model separates a current claim from a superseded one. In each case, the useful addition is a distinction that changes whether a state change is acceptable. The models make those distinctions available for analysis; the later Alignment chapter examines the mechanisms that enforce them.

% <!-- point: the-lease-model-omits-the-processing-topology | The model omits the processing topology: ownership, generation, lifetime, and reclamation are the consequential properties, and the implementation later changed how the job's work is partitioned without invalidating the model — the pointers to the code changed, not the model. -->

Notice what this model omits: the processing topology. Ownership, generation, lifetime, and reclamation matter to the question of who controls a job; how the job\textquotesingle s work is partitioned does not. DocAble\textquotesingle s implementation changed over time to improve performance while preserving those ownership semantics. We changed the pointers to the code, but not the model. The ownership properties could therefore remain invariants across those changes. The inset shows how a relationship built into the model can become an invariant when an implementation stores its parts separately.

% <!-- index-def: invariant -->

% <!-- box-family: inset -->

% <!-- inset-domain: FORMAL METHODS -->

% <!-- box-weight: callout -->

% mage-box-prefix: Inset —

\begin{bookbox}{inset}{What is an invariant?}
\label{inset-invariant}
An \textbf{invariant} is a property required to hold over a declared domain of states, structures, or observations. In a behavioral state machine, the invariant may be a predicate that must hold in every reachable state; in a structural model, it may require every observed dependency to belong to a declared set; in a measurement model, it may require a quantity to remain within a hard capacity bound. State both the property \emph{and} the domain the checker claims to establish it over.

Take a behavioral illustration. A job carries two fields at once --- a \texttt{status} of \texttt{free}, \texttt{leased}, or \texttt{done}, and a \texttt{lease\_owner} that names the holding worker or is empty. At the model level, leased-iff-owner was a definition; an implementation that represents the two sides as separately written fields turns it into a property that can fail, which is exactly why you write the invariant down. "A job is leased if and only if it has an owner" is a predicate over the pair: \texttt{(status~==~"leased")~==~(lease\_owner~is~not~None)}. A \texttt{free} status with a lingering \texttt{lease\_owner} is a stale lease; a \texttt{leased} status with no owner is a lost lease. Separating the two writes can make either inconsistent state reachable, and both break the predicate. The predicate becomes mechanically meaningful when a checker evaluates it over the declared reachable state space:

\begin{verbatim}
STATES = product({"free", "leased", "done"}, {None, "worker-1", "worker-2"})
def invariant(s):                                    # leased iff owned
    return (s.status == "leased") == (s.lease_owner is not None)
assert all(invariant(s) for s in reachable_states())  # check every reachable state
\end{verbatim}
\end{bookbox}

% <!-- point: the-lease-answers-ownership-behavioral-and-measurement-questions-at-once | Read as ownership the lease says who controls a job; read as behavioral it traces a claimed-held-expired-reclaimed lifecycle; counted, it measures how much work is in flight — one representation, three questions. -->

A lease illustrates how one representation can answer several classes of question. Its owner identifies who currently controls a job. Its generation and lifetime distinguish current claims from stale ones and support reasoning about expiry and reclamation. Counting the current claims across jobs gives a measurement of how much work is in flight. These questions concern ownership, behavior, and capacity, yet they draw on the same represented facts.

% <!-- point: overlap-confirms-the-classes-sort-questions-not-data-structures | The overlap follows from the distinction introduced at the chapter's opening: model classes identify engineering questions, not mutually exclusive data structures — the observation the closing section builds on to connect the models. -->

The overlap follows from the distinction introduced at the opening of the chapter: model classes identify engineering questions, not mutually exclusive data structures. One representation can serve several questions when it preserves the facts each requires. \magevalue{sec:scenarios-and-the-joined-model} develops this idea further by connecting models through shared identities.

Ownership introduced permission implicitly: a stale claimant \emph{may not} clear a current lease. The next section examines decisions directly, including permissions and how to choose what should happen next.
